\subsection{Long-lived assets, intangibles and goodwill}

These entries explain how intangible assets and goodwill are recognised, measured and disclosed, compare the reporting of intangible assets by how they were obtained, and explain impairment and derecognition of property, plant and equipment and intangibles, their effects on statements and ratios, and the disclosures from which an analyst infers asset age and policy.

\begin{qfmethod}[Intangible assets and goodwill on the balance sheet: recognition, measurement after recognition, impairment testing and the disclosures an analyst uses]\label{meth:intangibles-goodwill-reporting}
\qffield{Statement} \emph{Intangible assets.} An intangible asset is recognised only if future economic benefits are probable and its cost is reliably measurable \citep[IAS~38 para.~21]{eu-2023-1803-ias38}. Internally generated goodwill, brands, mastheads, publishing titles and customer lists are never recognised, because their cost cannot be separated from developing the business (paras.~48 and~63--64). After recognition an entity chooses cost or, only where an active market exists, revaluation (paras.~72 and~75). A finite-lived asset is amortised from when it is available for use, with its life and method reviewed each year as estimates (paras.~97 and~104). An indefinite-lived asset, with no foreseeable limit to its cash flows, is not amortised but tested for impairment annually and whenever there is an indication (paras.~88 and~107--108). Disclosures distinguish internally generated from other intangibles by class: lives or rates, methods, gross amounts, accumulated amortisation, a reconciliation of movements, and the research and development expensed (paras.~118 and~126).

\emph{Goodwill} is recognised only in a business combination. It is the excess of the consideration (plus any non-controlling interest and previously held interest) over the acquisition-date fair value of identifiable net assets \citep[IFRS~3 para.~32]{eu-2023-1803-ifrs3}. Identifiable intangibles, those that are separable or arise from contractual or legal rights, are recognised apart from goodwill (para.~B31). Goodwill is allocated to cash-generating units no larger than an operating segment and tested at least annually \citep[IAS~36 paras.~10 and~80]{eu-2023-1803-ias36}. An impairment loss reduces goodwill first, then other assets pro rata (para.~104), and is never reversed (para.~124). For units with significant goodwill, entities disclose the basis of recoverable amount and the key assumptions (para.~134).

\emph{US GAAP} does not amortise goodwill or indefinite-lived intangibles \citep[paras.~16--18]{us-goodwill-intangibles-2001}. Since 2017 the goodwill impairment charge has been the excess of a reporting unit's carrying amount over its fair value, capped at the goodwill allocated to it, rather than an implied-fair-value calculation \citep[pp.~2--3]{us-goodwill-2017}. SEC filers apply this to tests in fiscal years beginning after 15~December 2019.

\qffield{Notation} Goodwill $=$ consideration $+$ NCI $+$ prior interest $-$ fair value of identifiable net assets. The net assets include deferred tax on fair-value uplifts, which itself increases goodwill.

\qffield{Conditions} IFRS as adopted in the EU in 2023. US Statements are cited as issued; later Codification amendments, such as private-company alternatives, were not reviewed.

\qffield{Practice} \emph{Original.} An acquirer pays 500 for a business with book net assets of 250. Fair-value uplifts are:
\begin{itemize}
\item PP\&E 30, depreciated over six years;
\item customer relationships 90, amortised over ten years;
\item a brand 60 with indefinite life;
\item in-process research and development 40.
\end{itemize}
Deferred tax at 25~percent on the uplifts is 55, so identifiable net assets are 415 and goodwill is 85. Post-acquisition depreciation and amortisation are 14 a year; the brand, the in-process project (until completion) and goodwill are tested for impairment instead.

\emph{Analysis.} Harbour's goodwill of 120 is 9.4~percent of 2025 total assets (\hyperref[ex:harbour-tools-statements]{case}). The analyst reads the impairment-test assumptions, such as the discount rate, growth and headroom, against the forecast. Lack of amortisation means goodwill reflects past purchase prices, not current value. Many analysts therefore also examine returns on tangible capital.

\qffield{Pitfalls} Treating unamortised goodwill as evidence of value. Comparing an acquisitive company with one that grew organically, whose equivalent intangibles were expensed. Ignoring the sensitivity disclosures that show how close a unit is to impairment.

\qffield{See also} \hyperref[meth:intangible-acquisition-routes]{purchased, internal and acquired intangibles}; \hyperref[meth:impairment-derecognition]{impairment}.

\end{qfmethod}

\begin{qfmethod}[Purchased, internally developed and acquired intangible assets: recognition and measurement compared, and the effects on profit, assets and ratios]\label{meth:intangible-acquisition-routes}
\qffield{Statement}
\begin{itemize}
\item \emph{Purchased separately}: recognised at cost when benefits are probable and the cost is reliable, then amortised over a finite life or impairment-tested if indefinite \citep[IAS~38 paras.~21, 97 and~107--108]{eu-2023-1803-ias38}. US GAAP treats separately acquired intangibles in the same way \citep[paras.~16--17]{us-goodwill-intangibles-2001}.
\item \emph{Internally developed}: research is expensed; development is capitalised only after six conditions are demonstrated, and earlier spending is never reinstated (IAS~38 paras.~54, 57 and~71). Internally generated goodwill, brands and customer lists are never recognised (paras.~48 and~63). US GAAP expenses research and development \citep[para.~12]{us-research-development-1974}, except software costs after technological feasibility \citep[paras.~3--4]{us-software-1985} and internal-use software under Subtopic~350-40 (\hyperref[meth:expense-recognition-capitalisation]{capitalisation}).
\item \emph{Acquired in a business combination}: identifiable intangibles, including in-process research and development, are recognised separately from goodwill at acquisition-date fair value, even if the acquiree never recognised them (IAS~38 paras.~33--34; \citep[IFRS~3 paras.~10, 18 and~B31]{eu-2023-1803-ifrs3}). US GAAP recognises acquired research and development intangibles regardless of alternative future use and treats them as indefinite-lived until the project is completed or abandoned \citep[para.~66(a), pp.~20--21]{us-business-combinations-2007}.
\end{itemize}

\emph{Consequence.} The same economic asset can be on one balance sheet and absent from another. A brand bought with a company is recognised at fair value; a brand built internally is expensed through advertising. A development project acquired is capitalised; the same project developed in-house under US GAAP is expensed.

\qffield{Notation} Not applicable.

\qffield{Conditions} IFRS as adopted in 2023; US Statements as issued (FAS~2 of 1974, FAS~86 of 1985, FAS~142 of 2001, FAS~141R of 2007); later Codification amendments, such as private-company alternatives, are not reflected.

\qffield{Practice} \emph{Original.} Three companies each spend 40 on the same technology.
\begin{itemize}
\item A pays 40 to buy a finished patent; it recognises 40 and amortises 8 a year over five years under both frameworks.
\item B develops it internally; 24 is incurred after the IAS~38 conditions are met. An IFRS reporter expenses 16 and capitalises 24, while a US reporter expenses all 40, so its first-year expense is 24 higher and its assets 24 lower.
\item C acquires a company whose identifiable intangibles include the project. It recognises the in-process project at fair value with no amortisation until completion, and any excess consideration becomes goodwill (\hyperref[meth:intangibles-goodwill-reporting]{goodwill}).
\end{itemize}
Once spending is steady, return on assets is therefore highest for the US internal developer, whose investment is expensed and whose asset base excludes it, even though the economics are identical. The \hyperref[meth:accounting-analysis-adjustments]{comparability adjustment} puts companies on one basis, for example by capitalising and amortising disclosed development spending for all of them, or expensing it for all.

\qffield{Pitfalls} Ranking acquisitive and organic companies on return on assets without adjusting for unrecognised internal intangibles. Assuming an IFRS reporter's capitalised development is comparable across companies; meeting the six conditions is a judgement. Treating in-process research and development as amortising from acquisition.

\qffield{See also} \hyperref[meth:intangibles-goodwill-reporting]{intangibles and goodwill}; \hyperref[meth:impairment-derecognition]{impairment}.

\end{qfmethod}

\begin{qfmethod}[Impairment and derecognition of property, plant and equipment and intangible assets: IFRS and US tests, reversals, gains and losses on disposal, effects on statements and ratios, and reading the disclosures]\label{meth:impairment-derecognition}
\qffield{Statement} \emph{IFRS.} At each reporting date an entity assesses indications of impairment, such as market value declines, adverse changes, higher interest rates, obsolescence and worse performance \citep[IAS~36 paras.~9 and~12]{eu-2023-1803-ias36}. It tests intangibles not yet available for use, indefinite-lived intangibles and goodwill annually (para.~10). Recoverable amount is the higher of fair value less costs of disposal and value in use (para.~18). Value in use reflects expected cash flows, their variability, the time value of money and the price of uncertainty (para.~30). An asset is written down to recoverable amount through profit or loss, or through the revaluation surplus for a revalued asset (paras.~59--60). Losses on assets other than goodwill are reversed when estimates change, up to the carrying amount that would have existed without the impairment (paras.~110, 114 and~117); goodwill impairments are never reversed (para.~124).

\emph{US GAAP.} A long-lived asset is impaired only if its carrying amount is \emph{not recoverable}, meaning it exceeds the undiscounted cash flows from use and disposal. The loss is then the excess of carrying amount over fair value \citep[para.~7]{us-impairment-2001}. The written-down amount becomes the new cost basis, and restoring the loss is prohibited (para.~15). Assets held for sale are carried at the lower of carrying amount and fair value less costs to sell and are not depreciated (para.~34).

\emph{Derecognition.} The gain or loss on disposal is net proceeds less carrying amount, recognised in profit or loss but not as revenue \citep[IAS~16 paras.~67--68 and~71]{eu-2023-1803-ias16}. The proceeds are an investing inflow (\hyperref[meth:statement-linked-cash-flows]{cash flows}).

\emph{Effects.} An impairment reduces assets and equity and lowers profit in the year of loss, with no cash effect. Afterwards depreciation falls and the asset base is smaller, so return on assets and asset turnover rise without any operating improvement.

\emph{Disclosures} \citep[IAS~16 paras.~73 and~79]{eu-2023-1803-ias16}: gross carrying amounts, accumulated depreciation and a reconciliation of movements by class, with methods, lives and rates; idle and fully depreciated assets and fair values are encouraged. Impairment losses and reversals by class, with the events and key assumptions behind them, are also disclosed \citep[IAS~36 paras.~126, 130 and~134]{eu-2023-1803-ias36}.

\qffield{Notation} Average age $\approx$ accumulated depreciation$/$annual depreciation. Average total life $\approx$ gross cost$/$depreciation. Remaining life $\approx$ net carrying amount$/$depreciation. Straight-line depreciation is assumed.

\qffield{Conditions} The age measures assume straight-line depreciation and a stable asset mix; impairments and acquisitions distort them. FAS~144 (2001) is cited as issued; the current Codification was not consulted.

\qffield{Practice} \emph{Original.} A cash-generating unit carried at 200 has undiscounted cash flows of 230, fair value less costs of disposal of 150, value in use of 170 and fair value of 165.
\begin{itemize}
\item IFRS records an impairment of 30 (to 170); US GAAP records none, because the undiscounted flows exceed 200.
\item With five years of remaining life and operating profit of 50 before this asset's depreciation, return on the unit's average assets rises from 5.56 to 10.46~percent after the write-down.
\item Next year recoverable amount recovers to 190. IFRS reverses 24, restoring the 160 that would have remained without the impairment, and no more.
\end{itemize}
Harbour sold equipment in 2024 that had cost 40 with 30 accumulated depreciation for 13, a gain of 3 in operating profit and an investing inflow of 13.

\emph{Disclosure analysis} (\hyperref[ex:harbour-tools-statements]{Harbour}, 2025): accumulated depreciation of 440 over depreciation of 66 gives an average age of 6.67 years. Gross cost of 1{,}040 over 66 gives a total life of 15.76 years, so about 9.09 years remain. Capital expenditure of 105 is 1.6 times depreciation, consistent with growth and renewal, not a run-down of the asset base.

\qffield{Pitfalls} Treating post-impairment margin and return gains as operating improvement. Comparing IFRS and US companies' impairment frequency without the different tests. Taking disposal gains into recurring profit.

\qffield{See also} \hyperref[meth:intangible-acquisition-routes]{intangible routes}; \hyperref[meth:earnings-management-presentation]{big baths}.

\end{qfmethod}
